% =============================================================================
%  capitulos/cap3_metodologia.tex
%  CAPÍTULO 3 — METODOLOGIA
%
%  Objetivo: descrever COMO o trabalho foi desenvolvido do ponto de vista
%  de processo e gestão — não o que foi construído (isso vai no Cap. 5).
%
%  Responde: Qual metodologia ágil foi adotada? Como as etapas foram
%  organizadas? Quais ferramentas de gestão foram usadas?
% =============================================================================

\chapter{Metodologia}
\label{cap:metodologia}

TODO: Parágrafo introdutório explicando a abordagem metodológica geral
adotada no projeto e como este capítulo está organizado.


% -----------------------------------------------------------------------------
%  3.1 METODOLOGIA ÁGIL ADOTADA
%  Descreva a metodologia escolhida (Scrum, Kanban, XP, híbrida...).
%  Explique os ritos adotados e justifique a escolha para este contexto.
% -----------------------------------------------------------------------------
\section{Metodologia Ágil}
\label{sec:metodologia-agil}

TODO: Descreva a metodologia ágil escolhida e justifique a escolha.

% Exemplo para Scrum — adapte ou substitua conforme sua metodologia:
%
% O desenvolvimento seguiu a metodologia Scrum \cite{beck2004}, organizado
% em sprints de duas semanas. Os ritos adotados foram:
% \begin{itemize}
%   \item \textbf{Sprint Planning}: realizado no início de cada sprint para
%         definir as tarefas a serem desenvolvidas no período;
%   \item \textbf{Daily Standup}: reuniões diárias de 15 minutos para
%         alinhamento de progresso e identificação de impedimentos;
%   \item \textbf{Sprint Review}: apresentação das funcionalidades entregues
%         ao final de cada sprint;
%   \item \textbf{Retrospective}: avaliação do processo para identificar
%         melhorias na sprint seguinte.
% \end{itemize}


% -----------------------------------------------------------------------------
%  3.2 FERRAMENTAS DE GESTÃO E COMUNICAÇÃO
%  Liste as ferramentas usadas para organizar o projeto.
%  Ex.: Git, GitHub, Trello, Jira, Notion, Discord, Figma...
% -----------------------------------------------------------------------------
\section{Ferramentas de Gestão e Comunicação}
\label{sec:ferramentas-gestao}

TODO: Descreva as ferramentas utilizadas para cada finalidade:

\begin{itemize}
  \item \textbf{Controle de versão}: TODO: Git + GitHub/GitLab — repositório,
        branches, pull requests...
  \item \textbf{Gerenciamento de tarefas}: TODO: Trello / Jira / Notion /
        GitHub Projects — como o backlog foi organizado...
  \item \textbf{Comunicação}: TODO: Discord / Slack / WhatsApp...
  \item \textbf{Prototipação}: TODO: Figma / draw.io / Balsamiq...
  \item \textbf{Documentação}: TODO: Notion / Confluence / Google Docs...
\end{itemize}


% -----------------------------------------------------------------------------
%  3.3 ETAPAS DE DESENVOLVIMENTO
%  Descreva a divisão cronológica do trabalho.
%  Cada etapa deve ter: o que foi feito, quando ocorreu e qual foi a entrega.
% -----------------------------------------------------------------------------
\section{Etapas de Desenvolvimento}
\label{sec:etapas}

O desenvolvimento foi organizado nas seguintes etapas:

\begin{enumerate}
  \item \textbf{Levantamento de Requisitos}:
        TODO: Descreva como os requisitos foram coletados (entrevistas,
        questionários, análise de sistemas similares...), quando ocorreu
        e quais foram os artefatos gerados.

  \item \textbf{Modelagem e Prototipação}:
        TODO: Descreva a criação dos diagramas UML, protótipos de
        interface (wireframes/mockups) e validação com usuários/orientador.

  \item \textbf{Configuração do Ambiente}:
        TODO: Configuração do repositório, ambiente de desenvolvimento,
        CI/CD (se houver), conteinerização com Docker (se houver)...

  \item \textbf{Implementação do Backend}:
        TODO: Desenvolvimento das APIs, regras de negócio, integração
        com banco de dados...

  \item \textbf{Implementação do Frontend}:
        TODO: Desenvolvimento das interfaces, integração com o backend...

  \item \textbf{Testes e Validação}:
        TODO: Execução dos testes (unitários, integração, usabilidade)
        e correção de bugs identificados.

  \item \textbf{Ajustes Finais e Entrega}:
        TODO: Refatorações, documentação final e deploy (se aplicável).
\end{enumerate}


% -----------------------------------------------------------------------------
%  3.4 CRONOGRAMA (opcional)
%  Inclua se a sua instituição exigir ou se quiser documentar o planejamento.
%  Opção 1: tabela simples abaixo.
%  Opção 2: exporte o cronograma de uma ferramenta (ex.: Gantt no Notion)
%           como imagem e use \includegraphics.
% -----------------------------------------------------------------------------
% \section{Cronograma}
% \label{sec:cronograma}
%
% \begin{table}[H]
%   \caption{Cronograma de Desenvolvimento}
%   \label{tab:cronograma}
%   \centering
%   \begin{tabularx}{\textwidth}{Xcccccc}
%     \toprule
%     \textbf{Etapa} & \textbf{Mês 1} & \textbf{Mês 2} & \textbf{Mês 3}
%                    & \textbf{Mês 4} & \textbf{Mês 5} & \textbf{Mês 6} \\
%     \midrule
%     Levantamento de Requisitos  & X &   &   &   &   &   \\
%     Modelagem e Prototipação    & X & X &   &   &   &   \\
%     Implementação Backend       &   & X & X & X &   &   \\
%     Implementação Frontend      &   &   & X & X &   &   \\
%     Testes e Validação          &   &   &   & X & X &   \\
%     Ajustes e Entrega Final     &   &   &   &   & X & X \\
%     \bottomrule
%   \end{tabularx}
%   \fonte{Elaborado pelo autor.}
% \end{table}
